\documentclass[12pt]{article}
\usepackage[ukrainian,shorthands=off]{babel}   % shorthands=off: символ " не активний (потрібно для міток "..." у TikZ всередині \zadtask)
\usepackage{fontspec}
% --- НАЛАШТУВАННЯ ШРИФТІВ ---
% Шрифти Schoolbook-*.otf лежать у корені проєкту. Шлях підбирається автоматично:
% ../ (компіляція з папки теми), ./ (корінь проєкту), ../../ (підпапки теми 28); інакше -- TeX Gyre Schola.
\newcommand{\nmtSchoolbook}[1]{\setmainfont{Schoolbook-Regular.otf}[
    Path = #1,
    BoldFont = Schoolbook-Bold.otf,
    ItalicFont = Schoolbook-Italic.otf,
    BoldItalicFont = Schoolbook-BoldItalic.otf
]}
\IfFileExists{../Schoolbook-Regular.otf}{\nmtSchoolbook{../}}{%
 \IfFileExists{./Schoolbook-Regular.otf}{\nmtSchoolbook{./}}{%
  \IfFileExists{../../Schoolbook-Regular.otf}{\nmtSchoolbook{../../}}{%
   \setmainfont{texgyreschola}[Extension=.otf, UprightFont=*-regular, BoldFont=*-bold, ItalicFont=*-italic, BoldItalicFont=*-bolditalic]}}}
\usepackage[italic]{mathastext}
\usepackage[a4paper,margin=1.8cm,bottom=2cm,top=2cm]{geometry}
\usepackage{amsmath,amssymb}
\usepackage{tikz}
\usepackage{pgfplots}
\pgfplotsset{compat=1.18}
\usepgfplotslibrary{fillbetween}
\usetikzlibrary{calc,patterns,angles,quotes,intersections,babel,3d,shapes.symbols,shapes.geometric,decorations.pathreplacing,shadings,arrows.meta,hobby}
\usepackage{xcolor,array,fancyhdr,enumitem,multicol}
\usepackage{colortbl,multirow,diagbox,needspace}
\usepackage{tcolorbox}
\tcbuselibrary{skins,breakable}
\usepackage[hidelinks]{hyperref}
\usepackage[firstpage=false, text=@pvtr2525, scale=0.6, color=gray!12]{draftwatermark}

% --- АВТОСТИСКАННЯ: рисунок/сітка, ширша за колонку, масштабується до ширини колонки ---
\newsavebox{\nmtfitbox}
\newenvironment{nmtfit}{\begin{lrbox}{\nmtfitbox}}{\end{lrbox}%
  \ifdim\wd\nmtfitbox>\linewidth \resizebox{\linewidth}{!}{\usebox{\nmtfitbox}}\else\usebox{\nmtfitbox}\fi}
\let\nmtIncludegraphicsOrig\includegraphics
\renewcommand{\includegraphics}[2][]{\begin{nmtfit}\nmtIncludegraphicsOrig[#1]{#2}\end{nmtfit}}


\definecolor{mainGreen}{RGB}{34, 120, 64}
\definecolor{yearOrange}{RGB}{220, 100, 30}
\definecolor{yearcolor}{RGB}{220, 100, 30}
\definecolor{titleBg}{RGB}{225, 240, 220}
\definecolor{instrBg}{RGB}{255, 240, 220}
\definecolor{instrBorder}{RGB}{220, 150, 80}
\definecolor{headerblue}{RGB}{0, 102, 204}   % використовується в рисунках старої бази

\setlength{\parindent}{0pt}
\emergencystretch=1.5em

\pagestyle{fancy}
\fancyhf{}
\renewcommand{\headrulewidth}{0.4pt}
\renewcommand{\footrulewidth}{0.4pt}
\fancyhead[C]{\small\color{gray!80} Показникові вирази, функція, рівняння і нерівності}
\fancyhead[R]{\small\color{mainGreen}\textbf{\href{https://t.me/pvtr2525}{@pvtr2525}}}
\fancyfoot[L]{\small\color{gray!80}\href{https://t.me/pvtr2525}{ТГ: @pvtr2525}}
\fancyfoot[C]{\small\color{gray!80}\thepage}
\fancyfoot[R]{\small\color{gray!80}НМТ 2022--2026}

% --- ТАБЛИЦІ ВІДПОВІДЕЙ ---
% ширина клітинки = (ширина рядка - відступи) / 5: на всю сторінку це ~3 см, у minipage поруч із рисунком таблиця стискається
\newlength{\nmtcellw}
\newcommand{\nmtSetCell}{\setlength{\nmtcellw}{\dimexpr(\linewidth-12\tabcolsep-6\arrayrulewidth)/5\relax}}
\newcommand{\answerTable}[5]{%
\par\nopagebreak\vspace{0.15cm}
\noindent\nmtSetCell
\begin{tabular}{|*{5}{>{\centering\arraybackslash}m{\nmtcellw}|}}
\hline
\rule[-0.15cm]{0pt}{0.6cm}\textbf{А} & \rule[-0.15cm]{0pt}{0.6cm}\textbf{Б} & \rule[-0.15cm]{0pt}{0.6cm}\textbf{В} & \rule[-0.15cm]{0pt}{0.6cm}\textbf{Г} & \rule[-0.15cm]{0pt}{0.6cm}\textbf{Д} \\
\hline
\rule[-0.2cm]{0pt}{0.75cm}#1 & \rule[-0.2cm]{0pt}{0.75cm}#2 & \rule[-0.2cm]{0pt}{0.75cm}#3 & \rule[-0.2cm]{0pt}{0.75cm}#4 & \rule[-0.2cm]{0pt}{0.75cm}#5 \\
\hline
\end{tabular}
\par\vspace{0.25cm}
}
\newcommand{\answerTableTall}[5]{%
\par\nopagebreak\vspace{0.15cm}
\noindent\nmtSetCell
\begin{tabular}{|*{5}{>{\centering\arraybackslash}m{\nmtcellw}|}}
\hline
\rule[-0.2cm]{0pt}{0.7cm}\textbf{А} & \rule[-0.2cm]{0pt}{0.7cm}\textbf{Б} & \rule[-0.2cm]{0pt}{0.7cm}\textbf{В} & \rule[-0.2cm]{0pt}{0.7cm}\textbf{Г} & \rule[-0.2cm]{0pt}{0.7cm}\textbf{Д} \\
\hline
\rule[-0.45cm]{0pt}{1.2cm}#1 & \rule[-0.45cm]{0pt}{1.2cm}#2 & \rule[-0.45cm]{0pt}{1.2cm}#3 & \rule[-0.45cm]{0pt}{1.2cm}#4 & \rule[-0.45cm]{0pt}{1.2cm}#5 \\
\hline
\end{tabular}
\par\vspace{0.25cm}
}
\newcommand{\instructionBox}[1]{%
\par\vspace{0.3cm}
\begin{tcolorbox}[colback=instrBg, colframe=instrBorder, boxrule=0.6pt, arc=2pt,
    left=8pt, right=8pt, top=4pt, bottom=4pt]
\centering\small\bfseries #1
\end{tcolorbox}
\par\vspace{0.15cm}
}
\newcommand{\sectionTitle}[1]{%
\par\vspace{0.4cm}
\begin{tcolorbox}[colback=titleBg, colframe=titleBg, boxrule=0pt, arc=8pt,
    halign=center, fontupper=\Large\bfseries\color{mainGreen}]
\hyphenpenalty=10000 \exhyphenpenalty=10000 #1
\end{tcolorbox}
\par\vspace{0.3cm}
}
\newcommand{\task}[2]{\noindent\textbf{#1.}\ \ #2\par\vspace{0.2cm}}
% --- НМТ-СТИЛЬ ПОЛЕ ВІДПОВІДІ ---
\newcommand{\nmtAnswerBox}{%
\par\nopagebreak\vspace{0.2cm}\noindent
Відповідь:\ \
\begin{tabular}{@{}*{4}{|>{\centering\arraybackslash}p{0.8cm}}|@{\hspace{0.1cm}\textbf{,}\hspace{0.1cm}}*{3}{|>{\centering\arraybackslash}p{0.8cm}}|@{}}
\hline
\rule[-0.2cm]{0pt}{0.7cm} & & & & & & \\
\hline
\end{tabular}
\par\vspace{0.3cm}
}
% --- СІТКА ДЛЯ МАТЧИНГУ ---
\newsavebox{\nmtgridbox}
\newcommand{\matchingGrid}{%
    \sbox{\nmtgridbox}{\matchingGridRaw}%
    \ifdim\wd\nmtgridbox>\linewidth \resizebox{\linewidth}{!}{\usebox{\nmtgridbox}}\else\usebox{\nmtgridbox}\fi}
\newcommand{\matchingGridRaw}{%
    \begingroup
    \renewcommand{\arraystretch}{1.15}
    \setlength{\tabcolsep}{4pt}
    \footnotesize
    \begin{tabular}{r|c|c|c|c|c|}
         \multicolumn{1}{c}{} & \multicolumn{1}{c}{\textbf{А}} & \multicolumn{1}{c}{\textbf{Б}} & \multicolumn{1}{c}{\textbf{В}} & \multicolumn{1}{c}{\textbf{Г}} & \multicolumn{1}{c}{\textbf{Д}} \\ \cline{2-6}
         \textbf{1} & \phantom{X} & \phantom{X} & \phantom{X} & \phantom{X} & \phantom{X} \\ \cline{2-6}
         \textbf{2} & & & & & \\ \cline{2-6}
         \textbf{3} & & & & & \\ \cline{2-6}
    \end{tabular}
    \endgroup
}
% --- ДОДАТКОВІ МАКРОСИ ЗБІРНИКА ---
\newcommand{\taskBlock}[1]{%
\par\vspace{0.45cm}
\noindent{\color{mainGreen}\rule{\linewidth}{0.8pt}}\par\vspace{0.12cm}
\noindent{\small\bfseries\color{mainGreen}#1}\par\vspace{0.25cm}
}
\newcounter{zad}
\newcommand{\zadtask}[1]{\stepcounter{zad}\noindent\textbf{\thezad.}\ \ #1\par\nopagebreak\vspace{0.2cm}}
\newcommand{\zadnum}{\stepcounter{zad}\textbf{\thezad.}\ \ }
\newcounter{ans}
\newcommand{\ansitem}[1]{\stepcounter{ans}\noindent\textbf{\theans.}\ #1\par\vspace{0.07cm}}
% мітка року: нерозривна, притиснута праворуч; якщо не вміщається -- праворуч на новому рядку
\newcommand{\nmtyear}[1]{\unskip\nobreak\finalhyphendemerits=0 \hfill\penalty100\hskip1em\hbox{}\nobreak\hfill{\small\color{yearOrange}\mbox{\rule{0pt}{2.3ex}(НМТ~#1)}}}
% --- МАКРОС КОРОТКОГО РОЗВ'ЯЗКУ ---
\newcommand{\solution}[1]{%
\par\vspace{0.12cm}
\begin{tcolorbox}[colback=titleBg!45, colframe=mainGreen!55, boxrule=0.4pt, arc=2pt,
    left=6pt, right=6pt, top=3pt, bottom=3pt, breakable]
\footnotesize\textbf{Короткий розв'язок.}\ #1
\end{tcolorbox}
\par\vspace{0.12cm}
}
% Розв'язки й відповіді (є до завдань НМТ-2026): \showsolutionstrue -- показати, \showsolutionsfalse -- сховати
\newif\ifshowsolutions
\showsolutionsfalse
% --- ЗАГОЛОВКИ З ПУНКТАМИ КЛІКАБЕЛЬНОГО ЗМІСТУ ---
\setcounter{tocdepth}{2}
\newcommand{\chapterTitle}[1]{%
\clearpage\phantomsection\addcontentsline{toc}{section}{#1}%
\sectionTitle{#1}}
\newcommand{\typeTitle}[1]{%
\needspace{6\baselineskip}%
\phantomsection\addcontentsline{toc}{subsection}{\quad #1}%
\par\vspace{0.3cm}
\noindent{\large\bfseries\color{yearOrange}#1}\par\vspace{0.08cm}
\noindent{\color{yearOrange}\rule{\linewidth}{0.4pt}}\par\nopagebreak\vspace{0.2cm}}
\raggedbottom
% усередині одного завдання сторінка не розривається: нерозривний вертикальний проміжок
% і нерозривна межа абзаців (умова ніколи не відділяється від варіантів, рисунка чи сітки)
\newcommand{\nbvspace}[1]{\nopagebreak\vspace{#1}}
\newcommand{\nmtnobreak}{\ifvmode\penalty10000 \fi}
% --- ЄДИНИЙ МАКЕТ ЗАВДАНЬ НА ВІДПОВІДНІСТЬ ---
% мітка в колонці сталої ширини, текст -- у parbox: продовження довгого рядка
% вирівнюється під текстом, а не під міткою; відступ між пунктами однаковий скрізь
\newlength{\nmtMatchLab}\setlength{\nmtMatchLab}{1.6em}
\newlength{\nmtMatchSep}\setlength{\nmtMatchSep}{0.28cm}
\newcommand{\matchHead}[1]{\par\noindent\textit{#1}\par\nopagebreak\vspace{0.2cm}}
\newcommand{\matchItem}[2]{\par\noindent\makebox[\nmtMatchLab][l]{\textbf{#1}}%
\parbox[t]{\dimexpr\linewidth-\nmtMatchLab\relax}{\raggedright #2}\par\nopagebreak\vspace{\nmtMatchSep}}
\newcommand{\ansTheme}[1]{\par\vspace{0.25cm}\noindent{\bfseries\color{mainGreen}#1}\par\vspace{0.12cm}}
\newcommand{\ansType}[1]{\par\vspace{0.1cm}\noindent{\itshape\color{yearOrange}#1}\par\vspace{0.08cm}}

% варіанти відповіді окремими рядками (довгі вирази не влазять у клітинки таблиці)
\newcommand{\answerRows}[5]{\par\nopagebreak\vspace{0.15cm}%
\matchItem{А}{#1}\matchItem{Б}{#2}\matchItem{В}{#3}\matchItem{Г}{#4}\matchItem{Д}{#5}}

\begin{document}
\begingroup\setcounter{zad}{0}
% --- макроси з файлів сесій НМТ-2026 (для рисунків) ---
\renewcommand{\headrulewidth}{0.4pt}
\renewcommand{\footrulewidth}{0.4pt}

% --- сумісність зі старою базою (діє, якщо тема не має власного визначення) ---
\providecommand{\trigStyle}[1]{#1}
\providecommand{\answerTableSmall}[5]{%
\begingroup\setlength{\tabcolsep}{2pt}\small\nmtSetCell
\begin{tabular}{|*{5}{>{\centering\arraybackslash}m{\nmtcellw}|}}
\hline
\rule[-0.2cm]{0pt}{0.6cm}\textbf{А} & \textbf{Б} & \textbf{В} & \textbf{Г} & \textbf{Д} \\
\hline
\rule[-0.4cm]{0pt}{0.9cm}#1 & \rule[-0.4cm]{0pt}{0.9cm}#2 & \rule[-0.4cm]{0pt}{0.9cm}#3 & \rule[-0.4cm]{0pt}{0.9cm}#4 & \rule[-0.4cm]{0pt}{0.9cm}#5 \\
\hline
\end{tabular}\endgroup}
\providecommand{\matchingLayout}[3]{%
    \noindent
    \begin{minipage}[t]{0.36\textwidth}\vspace{0pt}\raggedright
        #1
    \end{minipage}%
    \hfill
    \begin{minipage}[t]{0.43\textwidth}\vspace{0pt}\raggedright
        #2
    \end{minipage}%
    \hfill
    \begin{minipage}[t]{0.19\textwidth}
        \vspace{0pt}
        \begin{flushright}
        #3
        \end{flushright}
    \end{minipage}}
\providecommand{\matchingLayoutBottom}[2]{%
    \noindent
    \begin{minipage}[t]{0.48\textwidth}\vspace{0pt}\raggedright
        #1
    \end{minipage}%
    \hfill
    \begin{minipage}[t]{0.48\textwidth}\vspace{0pt}\raggedright
        #2
    \end{minipage}
    \vspace{0.5cm}
    \noindent
    \matchingGrid}
\providecommand{\answerListVertical}[5]{%
    \par\nopagebreak\vspace{0.2cm}
    \begin{itemize}[itemsep=0.4cm, leftmargin=1.5cm, labelsep=0.5cm]
        \item[\textbf{А}] #1
        \item[\textbf{Б}] #2
        \item[\textbf{В}] #3
        \item[\textbf{Г}] #4
        \item[\textbf{Д}] #5
    \end{itemize}
    \vspace{0.2cm}}
\setcounter{zad}{0}

\chapterTitle{Показникові вирази, функція, рівняння і нерівності}

\typeTitle{Показникові вирази і функція}

\begin{samepage}
\noindent\zadnum Графік функції $y = 3^x$ паралельно перенесли на 2 одиниці вниз уздовж осі $y$. Укажіть функцію, графік якої отримали в результаті цього перетворення. \nmtyear{2025}\par
\nopagebreak\nbvspace{0.3cm}

\nmtnobreak
\answerTableTall{$y = 3^{x-2}$}{$y = 3^x - 2$}{$y = 3^{x+2}$}{$y = 3^x + 2$}{$y = \displaystyle\frac{3^x}{2}$}
\par\penalty-20
\end{samepage}

\begin{samepage}
\noindent\zadnum Спростіть вираз $5^{1+x} \cdot 5^{1-x}$. \nmtyear{2025}\par
\nopagebreak\nbvspace{0.3cm}

\nmtnobreak
\answerTable{$25^{1-x^2}$}{$5$}{$25$}{$5^{1-x^2}$}{$10$}
\par\penalty-20
\end{samepage}

\begin{samepage}
\zadtask{До виразу $2^{3a}$ доберіть тотожно рівний йому вираз \mbox{(А--Д)}, якщо $a \neq 8$. \nmtyear{2023}}
\answerRows{$8^a$}{$8 + a$}{$8a$}{$8 - a$}{$\displaystyle\frac{a}{8}$}
\par\penalty-20
\end{samepage}

\begin{samepage}
\zadtask{Установіть відповідність між виразом $2^\pi$ і проміжком \mbox{(А--Д)}, якому належить значення цього виразу. \nmtyear{2023}}
\answerRows{$(-\infty; 0]$}{$(0; 2)$}{$[2; 4)$}{$[4; 8)$}{$[8; +\infty)$}
\par\penalty-20
\end{samepage}

\begin{samepage}
\zadtask{Доберіть закінчення \mbox{(А--Д)} до початку речення <<Якщо $\displaystyle\frac{2^m}{2^n} = 2^a$, то>> так, щоб утворилося правильне твердження, якщо $m$ і $n$ --- натуральні числа, $n > 1$, $m > 1$. \nmtyear{2023}}
\answerRows{$a = mn$.}{$a = 0$.}{$a = m - n$.}{$a = n$.}{$a = \displaystyle\frac{m}{n}$.}
\par\penalty-20
\end{samepage}

\begin{samepage}
\zadtask{Установіть відповідність між функцією $y = 2^x$ та її найменшим значенням \mbox{(А--Д)} на відрізку $[0; 4]$. \nmtyear{2023}}
\answerRows{$1$}{$2$}{$3$}{$4$}{$5$}
\par\penalty-20
\end{samepage}

\begin{samepage}
\zadtask{Доберіть закінчення \mbox{(А--Д)} до початку речення <<Якщо $3^m \cdot 9^n = 3^a$, то>> так, щоб утворилося правильне твердження. \nmtyear{2023}}
\answerRows{$a = 0$.}{$a = 27^{mn}$.}{$a = 3mn$.}{$a = 2mn$.}{$a = m + 2n$.}
\par\penalty-20
\end{samepage}

\begin{samepage}
\noindent\zadnum \begin{minipage}[t]{0.55\textwidth}
У прямокутній системі координат на площині зображено коло $x^2+y^2=4$. Установіть відповідність між функцією $y=4^x$ та кількістю (А – Д) спільних точок, які має графік цієї функції із заданим колом. \nmtyear{2023}
\end{minipage}
\hfill
\begin{minipage}[t]{0.4\textwidth}
    \nopagebreak\nbvspace{-0.3cm}
    \begin{flushright}
    \begin{nmtfit}\begin{tikzpicture}[scale=0.6]
        % Сітка та осі

        \draw[->, >=stealth, thick] (-3,0) -- (3,0) node[below] {$x$};
        \draw[->, >=stealth, thick] (0,-3) -- (0,3) node[left] {$y$};

        % Коло
        \draw[thick] (0,0) circle (2cm);

        % Підписи
        \node[below left] at (0,0) {$0$};
        \node[above right] at (1.4, 1.4) {$x^2+y^2=4$};
    \end{tikzpicture}\end{nmtfit}
    \end{flushright}
\end{minipage}

\nmtnobreak
\nopagebreak\nbvspace{0.3cm}

\nmtnobreak
\noindent
\par\nopagebreak\vspace{0.15cm}
\answerRows{жодної}{лише одна}{лише дві}{лише три}{більше за три}
\par\penalty-20
\end{samepage}

\begin{samepage}
\zadtask{Увідповідніть функцію $y = \left(\displaystyle\frac{1}{2}\right)^x$ із кількістю \mbox{(А--Д)} спільних точок її графіка з прямою $y = -x$. \nmtyear{2023}}
\answerRows{жодної}{одна}{дві}{три}{безліч}
\par\penalty-20
\end{samepage}

\begin{samepage}
\zadtask{Установіть відповідність між функцією $y = 5^{-x}$ та її найбільшим значенням \mbox{(А--Д)} на відрізку $[-1; 3]$. \nmtyear{2023}}
\answerRows{$1$}{$2$}{$3$}{$4$}{$5$}
\par\penalty-20
\end{samepage}

\begin{samepage}
\zadtask{Увідповідніть функцію $y = 2^{-x}$ із кількістю \mbox{(А--Д)} спільних точок її графіка з прямою $y = x + 3$. \nmtyear{2023}}
\answerRows{жодної}{одна}{дві}{три}{безліч}
\par\penalty-20
\end{samepage}

\begin{samepage}
\zadtask{Установіть відповідність між функцією $y = 3^x$ та її властивістю \mbox{(А--Д)}. \nmtyear{2023}}
\answerRows{проходить через початок координат}{не містить точок з від'ємною ординатою}{не має спільних точок з віссю $y$}{двічі перетинає графік прямої $y = 3$}{графік знаходиться лише в I чверті}
\par\penalty-20
\end{samepage}

\begin{samepage}
\zadtask{Доберіть закінчення \mbox{(А--Д)} до початку речення <<Якщо $3^n \cdot 3 = 3^a$, то>> так, щоб утворилося правильне твердження, якщо $n > 0$. \nmtyear{2023}}
\answerRows{$a = 3n$.}{$a = n + 1$.}{$a = n + 3$.}{$a = \displaystyle\frac{3}{n}$.}{$a = \displaystyle\frac{n}{3}$.}
\par\penalty-20
\end{samepage}

\begin{samepage}
\zadtask{Установіть відповідність між твердженням <<графік функції проходить через точку $(0; 1)$>> та функцією \mbox{(А--Д)}, для якої це твердження є правильним. \nmtyear{2023}}
\answerRows{$y = -\cos x$}{$y = 4^x$}{$y = 2\sqrt{x + 1}$}{$y = x^2 - 4x + 3$}{$y = x$}
\par\penalty-20
\end{samepage}

\begin{samepage}
\zadtask{Установіть відповідність між функцією $y = 2^x$ та її властивістю \mbox{(А--Д)}. \nmtyear{2023}}
\answerRows{є зростаючою на $(-\infty; +\infty)$}{графік має одну спільну точку із графіком $y = x - 2$}{графік має дві спільні точки із графіком $y = x - 2$}{є спадною на проміжку $(-\infty; -2]$}{є спадною на проміжку $(-\infty; 2]$}
\par\penalty-20
\end{samepage}

\begin{samepage}
\zadtask{Установіть відповідність між виразом $2^{-a}$ і множиною \mbox{(А--Д)}, якій належить значення цього виразу, якщо $a = 5$. \nmtyear{2023}}
\answerRows{ірраціональних додатних чисел}{натуральних чисел}{цілих від'ємних чисел}{раціональних нецілих чисел}{ірраціональних від'ємних чисел}
\par\penalty-20
\end{samepage}

\begin{samepage}
\zadtask{Доберіть закінчення \mbox{(А--Д)} до початку речення <<Функція $y = 3^{-x}$>> так, щоб утворилося правильне твердження. \nmtyear{2023}}
\answerRows{має точку локального максимуму.}{має точку локального мінімуму.}{є непарною.}{зростає на всій області визначення.}{набуває лише додатних значень.}
\par\penalty-20
\end{samepage}

\begin{samepage}
\zadtask{Установіть відповідність між функцією $y = 2^x$ та властивістю її графіка \mbox{(А--Д)}. \nmtyear{2023}}
\answerRows{проходить через точку $(1; 0)$}{не перетинає вісь $y$}{симетричний відносно осі $y$}{розміщений лише в I та II чвертях}{розміщений лише в I та IV чвертях}
\par\penalty-20
\end{samepage}

\begin{samepage}
\zadtask{Установіть відповідність між твердженням <<найменше значення функції на відрізку $[-1; 3]$ дорівнює $0{,}5$>> та функцією \mbox{(А--Д)}, для якої це твердження є правильним. \nmtyear{2024}}
\answerRows{$y = x^2 - 4$}{$y = \displaystyle\frac{1}{x - 4}$}{$y = 2^x$}{$y = 0{,}5^x$}{$y = \sqrt{x + 1}$}
\par\penalty-20
\end{samepage}

\begin{samepage}
\zadtask{Установіть відповідність між функцією $y = 4^x$ та її властивістю \mbox{(А--Д)}. \nmtyear{2024}}
\answerRows{є зростаючою на всій області визначення}{набуває від'ємного значення при $x = -2$}{є парною}{має три спільні точки з графіком функції $y = -x$}{графік функції розташований лише в I та IV координатній чверті}
\par\penalty-20
\end{samepage}

\begin{samepage}
\zadtask{Установіть відповідність між виразом $\left(\displaystyle\frac{1}{2}\right)^{\log_2 \pi}$ та проміжком \mbox{(А--Д)}, якому належить значення цього виразу. \nmtyear{2024}}
\answerRows{$[-5; -2)$}{$[-2; 0)$}{$[0; 1)$}{$[1; 2)$}{$[2; 5)$}
\par\penalty-20
\end{samepage}

\begin{samepage}
\zadtask{Установіть відповідність між функцією $y = 4^x$ та властивістю її графіка \mbox{(А--Д)}. \nmtyear{2024}}
\answerRows{розташований лише в першій чверті}{має спільну точку з віссю $x$}{проходить через точку $(0; 1)$}{не перетинає вісь $y$}{симетричний відносно осі $y$}
\par\penalty-20
\end{samepage}

\begin{samepage}
\zadtask{Доберіть закінчення \mbox{(А--Д)} до початку речення <<Графік функції $y = 3^x$ має лише одну спільну точку з>> так, щоб утворилося правильне твердження. \nmtyear{2024}}
\answerRows{віссю $x$.}{віссю $y$.}{прямою $y = x - 4$.}{прямою $y = -4$.}{прямою $y = -2$.}
\par\penalty-20
\end{samepage}

\begin{samepage}
\zadtask{Доберіть закінчення \mbox{(А--Д)} до початку речення <<Графік функції $y = 2^x$>> так, щоб утворилося правильне твердження. \nmtyear{2024}}
\answerRows{симетричний відносно осі абсцис.}{симетричний відносно осі ординат.}{симетричний відносно початку координат.}{не перетинає вісь абсцис.}{не перетинає вісь ординат.}
\par\penalty-20
\end{samepage}

\begin{samepage}
% рік визначено за сусідніми завданнями (у старому файлі тег року відсутній)
\noindent\zadnum Установіть відповідність між твердженням <<не має спільних точок з функцією $y=\left(\dfrac{2}{3}\right)^x$>> та прямою, зображеною на рисунку \mbox{(А--Д)}, для якої це твердження є правильним. \nmtyear{2024}

\nmtnobreak
\nopagebreak\nbvspace{0.3cm}

\nmtnobreak
\noindent
%
\hfill

\nmtnobreak
\nopagebreak\nbvspace{0.5cm}

\nmtnobreak
% Таблиця з графіками
\begin{center}
\begingroup
\setlength{\tabcolsep}{3pt}
\begin{tabular}{|*{5}{c|}}
\hline
\textbf{А} & \textbf{Б} & \textbf{В} & \textbf{Г} & \textbf{Д} \\
\hline
\begin{nmtfit}\begin{tikzpicture}[scale=0.35]
    \draw[->, >=stealth] (-1.5,0) -- (3.5,0) node[below] {$x$};
    \draw[->, >=stealth] (0,-3) -- (0,2) node[left] {$y$};
    \node[below left] at (0,0) {\tiny $0$};
    \node[below] at (2,0) {\tiny $2$};
    \node[left] at (0,-2) {\tiny $-2$};
    \draw[thick] (-0.5,-2.5) -- (2.5,0.5); % y = x - 2
\end{tikzpicture}\end{nmtfit} &
\begin{nmtfit}\begin{tikzpicture}[scale=0.35]
    \draw[->, >=stealth] (-2.5,0) -- (2.5,0) node[below] {$x$};
    \draw[->, >=stealth] (0,-2) -- (0,3) node[left] {$y$};
    \node[below left] at (0,0) {\tiny $0$};
    \node[above right] at (0,1) {\tiny $1$};
    \draw (0.1,1) -- (0,1) -- (0.2,1.2) -- (0.2,1); % прямий кут
    \draw[thick] (-2.5,1) -- (2.5,1); % y = 1
\end{tikzpicture}\end{nmtfit} &
\begin{nmtfit}\begin{tikzpicture}[scale=0.35]
    \draw[->, >=stealth] (-2.5,0) -- (2.5,0) node[below] {$x$};
    \draw[->, >=stealth] (0,-2) -- (0,3) node[left] {$y$};
    \node[below right] at (0,0) {\tiny $0$};
    \node[below left] at (-1,0) {\tiny $-1$};
    \draw (-1,0) -- (-0.8,0) -- (-0.8,0.2) -- (-1,0.2); % прямий кут
    \draw[thick] (-1,-2) -- (-1,3); % x = -1
\end{tikzpicture}\end{nmtfit} &
\begin{nmtfit}\begin{tikzpicture}[scale=0.35]
    \draw[->, >=stealth] (-3.5,0) -- (2.5,0) node[below] {$x$};
    \draw[->, >=stealth] (0,-3) -- (0,2) node[left] {$y$};
    \node[below right] at (0,0) {\tiny $0$};
    \node[below] at (-2,0) {\tiny $-2$};
    \node[left] at (0,-2) {\tiny $-2$};
    \draw[thick] (-2.5,0.5) -- (0.5,-2.5); % y = -x - 2
\end{tikzpicture}\end{nmtfit} &
\begin{nmtfit}\begin{tikzpicture}[scale=0.35]
    \draw[->, >=stealth] (-1.5,0) -- (4,0) node[below] {$x$};
    \draw[->, >=stealth] (0,-3) -- (0,3) node[left] {$y$};
    \node[below left] at (0,0) {\tiny $0$};
    \node[below] at (3,0) {\tiny $3$};
    \node[left] at (0,-2) {\tiny $-2$};
    \draw[thick] (-1,-2.66) -- (4,0.66); % y = 2/3 x - 2
\end{tikzpicture}\end{nmtfit} \\
\hline
\end{tabular}
\endgroup
\end{center}
\par\penalty-20\vspace{0.25cm}
\end{samepage}

\begin{samepage}
\zadtask{Узгодьте функцію $y = 2^x$ із кількістю \mbox{(А--Д)} спільних точок її графіка з прямою $y = -x$. \nmtyear{2025}}
\answerRows{жодної}{одна}{дві}{три}{безліч}
\par\penalty-20
\end{samepage}

\begin{samepage}
\zadtask{Узгодьте вираз $4^6 : 4^{-m}$ із значенням $m$ \mbox{(А--Д)}, за якого значення цього виразу дорівнює 1. \nmtyear{2025}}
\answerRows{$\displaystyle\frac{1}{4}$}{$6$}{$4$}{$-6$}{$8$}
\par\penalty-20
\end{samepage}

\begin{samepage}
\zadtask{Доберіть закінчення \mbox{(А--Д)} до початку речення <<Якщо $\sqrt{2^m} = 2^n$, то>> так, щоб утворилося правильне твердження, якщо $n \neq 0$. \nmtyear{2025}}
\answerRows{$n = 2m$.}{$n = m^2$.}{$n = m^2 + 2m$.}{$n = \sqrt{m}$.}{$n = \displaystyle\frac{m}{2}$.}
\par\penalty-20
\end{samepage}

\begin{samepage}
% Відповідь: Г
\zadtask{Доберіть закінчення \mbox{(А--Д)} до початку речення <<Графік функції $y = 4^x$>> так, щоб утворилося правильне твердження. \nmtyear{2026}}
\answerRows{симетричний відносно початку координат.}{симетричний відносно осі $y$.}{перетинає лише вісь $x$.}{перетинає лише вісь $y$.}{перетинає графік функції $y = x$ лише в одній точці.}
\par\penalty-20
\end{samepage}

\begin{samepage}
% Відповідь: Д
\zadtask{Узгодьте вираз $\sqrt{\left(2^a\right)^3}$ із його значенням \mbox{(А--Д)}, якщо $a = \dfrac{2}{3}$. \nmtyear{2026}}
\answerRows{$4$}{$-1$}{$1$}{$\dfrac{20}{3}$}{$2$}
\par\penalty-20
\end{samepage}

\begin{samepage}
\zadtask{Доберіть до функції $y = 0{,}2^x - 1$ її властивість \mbox{(А--Д)}. \nmtyear{2026}}
\answerRows{є зростаючою}{є спадною}{не має нулів}{має лише два нулі}{є парною}
\par\penalty-20
\end{samepage}

\begin{samepage}
\zadtask{Доберіть до функції $y = 3^x$ її властивість \mbox{(А--Д)}. \nmtyear{2026}}
\answerRows{областю значень є проміжок $[3; +\infty)$}{областю визначення є проміжок $[3; +\infty)$}{областю визначення є проміжок $[-3; +\infty)$}{зростає на множині всіх дійсних чисел}{є непарною}
\par\penalty-20
\end{samepage}

\begin{samepage}
% Відповідь: Б
\zadtask{Доберіть закінчення \mbox{(А--Д)} до початку речення <<Графіки функцій $y = 2^x$ та $y = 0{,}5^x$ перетинаються в точці>> так, щоб утворилося правильне твердження. \nmtyear{2026}}
\answerRows{$(1; 0)$.}{$(0; 1)$.}{$(2; 0)$.}{$(-2; 0)$.}{$(-1; 0)$.}
\par\penalty-20
\end{samepage}

\begin{samepage}
% Відповідь: В
\zadtask{Доберіть до функції $y = -5^x$ властивість її графіка \mbox{(А--Д)}. \nmtyear{2026}}
\answerRows{симетричний відносно початку координат}{симетричний відносно осі $y$}{розміщений лише в ІІІ та IV координатних чвертях}{розміщений лише в І координатній чверті}{проходить через точку $(0; 1)$}
\par\penalty-20
\end{samepage}

\typeTitle{Показникові рівняння}

\begin{samepage}
\noindent\zadnum Розв'яжіть систему рівнянь $\begin{cases} 2^{3x-y} = 2^7, \\ 2x + y = 3. \end{cases}$ Якщо $(x_0; y_0)$ --- розв'язок цієї системи, то $x_0 \cdot y_0 =$ \nmtyear{2023}

\nmtnobreak
\answerTable{$2$}{$1$}{$-3$}{$-1$}{$-2$}
\par\penalty-20
\end{samepage}

\begin{samepage}
\noindent\zadnum Укажіть проміжок, якому належить корінь рівняння $\left(\displaystyle\frac{1}{3}\right)^{2x-1} = 9$. \nmtyear{2023}

\nmtnobreak
\answerTable{$[2; +\infty)$}{$(-2; 0]$}{$(-\infty; -2]$}{$[1; 2)$}{$(0; 1]$}
\par\penalty-20
\end{samepage}

\begin{samepage}
\noindent\zadnum Розв'яжіть рівняння $3^x = 81 \cdot 9^x$. \nmtyear{2023}

\nmtnobreak
\answerTable{$4$}{$-4$}{$-2$}{$-3$}{$3$}
\par\penalty-20
\end{samepage}

\begin{samepage}
% === ЗАВДАННЯ ===
\noindent\zadnum Якому проміжку належить корінь рівняння $2^x = \displaystyle\frac{1}{8}$? \nmtyear{2023}

\nmtnobreak
\answerTable{$(2; 4]$}{$(-6; -4]$}{$(-2; 0]$}{$(0; 2]$}{$(-4; -2]$}

\nmtnobreak
% === НМТ 2024 ===
\par\penalty-20\vspace{0.25cm}
\end{samepage}

\begin{samepage}
\noindent\zadnum Укажіть проміжок, якому належить корінь рівняння $4^x \cdot 5^x = \displaystyle\frac{1}{400}$. \nmtyear{2024}

\nmtnobreak
\answerTable{$[2; +\infty)$}{$[-0{,}5; 2)$}{$[-2; -0{,}5)$}{$[-10; -2)$}{$(-\infty; -10)$}
\par\penalty-20
\end{samepage}

\begin{samepage}
\noindent\zadnum Укажіть проміжок, якому належить корінь рівняння $5^{3x} \cdot 5^{-2x} = \displaystyle\frac{1}{5}$. \nmtyear{2024}

\nmtnobreak
\answerTable{$(-\infty; -1]$}{$(-1; -0{,}5]$}{$(-0{,}5; 0]$}{$(0{,}5; +\infty)$}{$(0; 0{,}5]$}
\par\penalty-20
\end{samepage}

\begin{samepage}
\noindent\zadnum Укажіть корінь рівняння $2^x = 2 \cdot 16$. \nmtyear{2024}

\nmtnobreak
\answerTable{$7$}{$3$}{$6$}{$4$}{$5$}
\par\penalty-20
\end{samepage}

\begin{samepage}
\noindent\zadnum Укажіть проміжок, якому належить корінь рівняння $\left(\displaystyle\frac{1}{3}\right)^{x-7} = \sqrt{3}$. \nmtyear{2024}

\nmtnobreak
\answerTable{$(7; 10]$}{$(-4; 4)$}{$(-6; -4]$}{$(4; 7]$}{$[-7; -6]$}
\par\penalty-20
\end{samepage}

\begin{samepage}
\noindent\zadnum Укажіть проміжок, якому належить \textit{менший} корінь рівняння $3^{x^2} = 81$. \nmtyear{2024}

\nmtnobreak
\answerTable{$[-2; 0)$}{$(2; +\infty)$}{$[-3; -2)$}{$[0; 2)$}{$(-\infty; -3)$}
\par\penalty-20
\end{samepage}

\begin{samepage}
\noindent\zadnum Розв'яжіть систему рівнянь $\begin{cases} 2^x + 2y = 0, \\ 2^x - 4y = 6. \end{cases}$ \nmtyear{2024}

\nmtnobreak
\answerTable{$(-3; 3)$}{$(1; -1)$}{$(3; -3)$}{$(-1; 1)$}{$(2; -2)$}
\par\penalty-20
\end{samepage}

\begin{samepage}
\noindent\zadnum Розв'яжіть рівняння $10^x = 0{,}1$. \nmtyear{2024}

\nmtnobreak
\answerTable{$1$}{$-1$}{$0$}{$0{,}01$}{$-9{,}9$}

\nmtnobreak
% === НМТ 2025 ===
\par\penalty-20\vspace{0.25cm}
\end{samepage}

\begin{samepage}
\noindent\zadnum Розв'яжіть систему рівнянь $\begin{cases} 2^{y+2x} = 64, \\ 2x - y = 12. \end{cases}$ Якщо $(x_0; y_0)$ --- розв'язок системи, то $x_0 + y_0 =$ \nmtyear{2025}

\nmtnobreak
\answerTable{$4{,}5$}{$1{,}5$}{$7{,}5$}{$-3$}{$9$}
\par\penalty-20
\end{samepage}

\begin{samepage}
\noindent\zadnum Розв'яжіть рівняння $3^{x+1} = 81$. \nmtyear{2025}

\nmtnobreak
\answerTable{$5$}{$3$}{$-3$}{$4$}{$2$}
\par\penalty-20
\end{samepage}

\begin{samepage}
\noindent\zadnum Укажіть проміжок, якому належить корінь рівняння $\left(\displaystyle\frac{3}{8}\right)^x = \displaystyle\frac{8}{3}$. \nmtyear{2025}

\nmtnobreak
\answerTable{$(-2; 0)$}{$(2; 4)$}{$(-6; -4)$}{$(-4; -2)$}{$(0; 2)$}
\par\penalty-20
\end{samepage}

\begin{samepage}
\noindent\zadnum Укажіть проміжок, якому належить корінь рівняння $\left(\displaystyle\frac{1}{3}\right)^{x-7} = \sqrt{3}$. \nmtyear{2025}

\nmtnobreak
\answerTable{$(-7; -6)$}{$(-6; 0)$}{$(0; 1)$}{$(6; 7)$}{$(1; 6)$}
\par\penalty-20
\end{samepage}

\begin{samepage}
% НМТ 2026, сесія 23.05, завдання №15 --- основна тема: Системи рівнянь
\zadtask{Розв'яжіть систему рівнянь $\begin{cases} 5^x = \sqrt{5}, \\ \dfrac{y + 1}{2x + 3} = \dfrac{1 - 4x}{2}. \end{cases}$ Якщо $(x_0;\, y_0)$~--- розв'язок системи, то $x_0 + y_0 = $ \nmtyear{2026}}
\answerTable{$-2{,}5$}{$-1{,}5$}{$2{,}5$}{$-3{,}5$}{$1{,}5$}
% Відповідь: А
\ifshowsolutions\solution{З першого рівняння $5^x=5^{1/2}$, тож $x_0=0{,}5$. Підставимо: $\dfrac{y+1}{2\cdot 0{,}5+3}=\dfrac{1-4\cdot 0{,}5}{2}$, тобто $\dfrac{y+1}{4}=-\dfrac{1}{2}$, звідки $y+1=-2$, $y_0=-3$. Тоді $x_0+y_0=0{,}5-3=-2{,}5$. Відповідь: \textbf{А}.}\fi
\par\penalty-20
\end{samepage}

\begin{samepage}
% НМТ 2026, сесія 1.06, завдання №1
\zadtask{Позначте число, що є коренем рівняння $3 \cdot 3^x = 27$. \nmtyear{2026}}
\answerTable{$2$}{$3$}{$4$}{$9$}{$27$}
% Відповідь: А
\ifshowsolutions\solution{$3\cdot 3^x=27 \Rightarrow 3^{x+1}=3^3 \Rightarrow x+1=3 \Rightarrow x=2$. Відповідь: \textbf{А}.}\fi
\par\penalty-20
\end{samepage}

\begin{samepage}
% НМТ 2026, сесія 3.06, завдання №11
\zadtask{Якому проміжку належить корінь рівняння $2^{x+1} - 2^x = \dfrac{1}{8}$? \nmtyear{2026}}
\answerTable{$(-\infty;\, -5)$}{$[-5;\, -3)$}{$[-3;\, 0)$}{$[0;\, 5)$}{$[5;\, +\infty)$}
% Відповідь: В
\ifshowsolutions\solution{$2^{x+1} - 2^x = 2^x(2-1) = 2^x = \dfrac18 = 2^{-3}$, звідки $x = -3 \in [-3;\, 0)$. Відповідь: \textbf{В}.}\fi
\par\penalty-20
\end{samepage}

\begin{samepage}
% НМТ 2026, сесія 5.06, завдання №13
\zadtask{Якому проміжку належить корінь рівняння $5^x \cdot 20^x = 1000$? \nmtyear{2026}}
\answerTable{$(-\infty;\, 0)$}{$[0;\, 1{,}5)$}{$[1{,}5;\, 2)$}{$[2;\, 4)$}{$[4;\, +\infty)$}
% Відповідь: В
\ifshowsolutions\solution{$5^x\cdot 20^x=(5\cdot 20)^x=100^x=10^{2x}$, а $1000=10^3$. Отже $10^{2x}=10^3\Rightarrow 2x=3\Rightarrow x=1{,}5$, що належить $[1{,}5;\,2)$. Відповідь: \textbf{В}.}\fi
\par\penalty-20
\end{samepage}

\begin{samepage}
% НМТ 2026, сесія 8.06, завдання №7
\zadtask{Якому проміжку належить корінь рівняння $2^{2x-5} = 8$? \nmtyear{2026}}
\answerTable{$(-\infty; -2)$}{$[0; 2)$}{$[-2; 0)$}{$[2; 4)$}{$[4; +\infty)$}
\par\penalty-20
\end{samepage}

\begin{samepage}
% НМТ 2026, сесія 9.06, завдання №7
\zadtask{Розв'яжіть рівняння $3^{x} = 81 \cdot 9^{x}$. \nmtyear{2026}}
\answerTable{$-2$}{$-4$}{$-3$}{$3$}{$4$}
\par\penalty-20
\end{samepage}

\begin{samepage}
% НМТ 2026, сесія 10.06, завдання №7
\zadtask{Позначте проміжок, якому належить корінь рівняння $\left(\dfrac{1}{3}\right)^{2x-1} = 9$. \nmtyear{2026}}
\answerTable{$(-3; -2)$}{$(-2; -1)$}{$(-1; 0)$}{$(0; 1)$}{$(1; 2)$}
\par\penalty-20
\end{samepage}

\begin{samepage}
% НМТ 2026, сесія 17.06, завдання №7
\zadtask{Розв'яжіть рівняння $3^x = 9^{x+2}$. \nmtyear{2026}}
\answerTable{$-4$}{$4$}{$0$}{$-2$}{$2$}
% Відповідь: А
\par\penalty-20
\end{samepage}

\begin{samepage}
% НМТ 2026, сесія 18.06, завдання №15 --- основна тема: Системи рівнянь
\zadtask{Розв'яжіть систему рівнянь $\begin{cases}2^{3x - y} = 2^7,\\ 2x + y = 3.\end{cases}$ Якщо $(x_0;\, y_0)$~--- розв'язок цієї системи, то $x_0 \cdot y_0 = $ \nmtyear{2026}}
\answerTable{$2$}{$-2$}{$-1$}{$-3$}{$1$}
% Відповідь: Б
\par\penalty-20
\end{samepage}

\begin{samepage}
% НМТ 2026, сесія 19.06, завдання №15
\zadtask{Визначте проміжок, якому належить корінь рівняння $5^{6x + 2} = \dfrac{1}{25}$. \nmtyear{2026}}
\answerTable{$[-2; -1)$}{$[-1; 0)$}{$[0; 1)$}{$[1; 2)$}{$[2; +\infty)$}
% Відповідь: Б
\par\penalty-20
\end{samepage}

\typeTitle{Показникові нерівності}

\begin{samepage}
\noindent\zadnum Яке з наведених чисел є розв'язком нерівності $5^{-x} > 25$? \nmtyear{2023}

\nmtnobreak
\answerTable{$0$}{$3$}{$-1$}{$-3$}{$1$}
\par\penalty-20
\end{samepage}

\begin{samepage}
\noindent\zadnum Розв'яжіть нерівність $\left(\displaystyle\frac{1}{5}\right)^x > \displaystyle\frac{1}{25}$. \nmtyear{2023}

\nmtnobreak
\answerTable{$(5; +\infty)$}{$(-\infty; 5)$}{$(2; +\infty)$}{$(0; 2)$}{$(-\infty; 2)$}
\par\penalty-20
\end{samepage}

\begin{samepage}
\noindent\zadnum Розв'яжіть нерівність $8^x > \displaystyle\frac{1}{64}$. \nmtyear{2023}

\nmtnobreak
\answerTable{$(-\infty; -2)$}{$(0; 2)$}{$(2; +\infty)$}{$(-2; +\infty)$}{$(-\infty; 2)$}

\nmtnobreak
% === НМТ 2024 ===
\par\penalty-20\vspace{0.25cm}
\end{samepage}

\begin{samepage}
\noindent\zadnum Розв'яжіть нерівність $\left(\displaystyle\frac{1}{3}\right)^{2x-1} \leqslant 27^{-1}$. \nmtyear{2024}

\nmtnobreak
\answerTable{$[2; +\infty)$}{$[-1; +\infty)$}{$(-\infty; 4]$}{$(-\infty; -1]$}{$(-\infty; 2]$}
\par\penalty-20
\end{samepage}

\begin{samepage}
\noindent\zadnum Розв'яжіть систему нерівностей $\begin{cases} 5^x < 25, \\ 2 - x < 8. \end{cases}$ \nmtyear{2024}

\nmtnobreak
\answerTable{$(-6; 5)$}{$(2; 6)$}{$(-\infty; -6)$}{$(-6; 2)$}{$(2; +\infty)$}
\par\penalty-20
\end{samepage}

\begin{samepage}
\noindent\zadnum Розв'яжіть систему нерівностей $\begin{cases} x^2 + 9 \geqslant 0, \\ 2^x > \displaystyle\frac{1}{16}. \end{cases}$ \nmtyear{2024}

\nmtnobreak
\answerTable{$[-3; +\infty)$}{$(-4; +\infty)$}{$(-4; -3]$}{$(-4; -3] \cup [3; +\infty)$}{$\varnothing$}
\par\penalty-20
\end{samepage}

\begin{samepage}
\noindent\zadnum Розв'яжіть нерівність $(0{,}1)^{x+2} < 0{,}1$. \nmtyear{2024}

\nmtnobreak
\answerTable{$(-1; +\infty)$}{$(-\infty; 3)$}{$(-\infty; -1)$}{$(-2; +\infty)$}{$(3; +\infty)$}
\par\penalty-20
\end{samepage}

\begin{samepage}
\noindent\zadnum Розв'яжіть нерівність $5^x \leqslant 1$. \nmtyear{2024}

\nmtnobreak
\answerTableTall{$[1; +\infty)$}{$(-\infty; 1]$}{$\left(-\infty; \displaystyle\frac{1}{5}\right)$}{$(-\infty; 0]$}{$[0; +\infty)$}

\nmtnobreak
% === НМТ 2025 ===
\par\penalty-20\vspace{0.25cm}
\end{samepage}

\begin{samepage}
\noindent\zadnum Розв'яжіть нерівність $7^x \leqslant \displaystyle\frac{1}{49}$. \nmtyear{2025}

\nmtnobreak
\answerTableTall{$\left(-\infty; \displaystyle\frac{1}{243}\right]$}{$[-2; +\infty)$}{$\left[\displaystyle\frac{1}{7}; +\infty\right)$}{$(-\infty; -2]$}{$\left(-\infty; \displaystyle\frac{1}{2}\right]$}
\par\penalty-20
\end{samepage}

\begin{samepage}
\noindent\zadnum Розв'яжіть нерівність $\left(\displaystyle\frac{1}{4}\right)^{2x-6} > \displaystyle\frac{1}{16}$. \nmtyear{2025}

\nmtnobreak
\answerTable{$(-\infty; -2)$}{$(-\infty; 2)$}{$(-2; +\infty)$}{$(4; +\infty)$}{$(-\infty; 4)$}
\par\penalty-20
\end{samepage}

\begin{samepage}
\noindent\zadnum Розв'яжіть нерівність $3^x \cdot 5^x \leqslant \displaystyle\frac{1}{15}$. \nmtyear{2025}

\nmtnobreak
\answerTableTall{$(-\infty; 1]$}{$[1; +\infty)$}{$\left(-\infty; \displaystyle\frac{1}{15}\right]$}{$[-1; +\infty)$}{$(-\infty; -1]$}
\par\penalty-20
\end{samepage}

\begin{samepage}
\noindent\zadnum Розв'яжіть нерівність $2^{3-x} > \displaystyle\frac{1}{4}$. \nmtyear{2025}

\nmtnobreak
\answerTable{$(-\infty; 5)$}{$(-5; +\infty)$}{$(-1; +\infty)$}{$(5; +\infty)$}{$(-\infty; -1)$}
\par\penalty-20
\end{samepage}

\begin{samepage}
\noindent\zadnum Розв'яжіть нерівність $10^{x-2} > 1$. \nmtyear{2025}

\nmtnobreak
\answerTable{$(3; +\infty)$}{$(-2; +\infty)$}{$(-\infty; 2)$}{$(2{,}1; +\infty)$}{$(2; +\infty)$}
\par\penalty-20
\end{samepage}

\begin{samepage}
% НМТ 2026, сесія 4.06, завдання №15
\zadtask{Розв'яжіть нерівність $3^{x^2 - 2x + 1} \leqslant 1$. \nmtyear{2026}}
\answerTable{$1$}{$(-\infty;\, -1]$}{$-1$}{$[1;\, +\infty)$}{$[-1;\, +\infty)$}
% Відповідь: А
\ifshowsolutions\solution{Оскільки $1=3^0$ і основа $3>1$, нерівність рівносильна $x^2-2x+1\leqslant0$, тобто $(x-1)^2\leqslant0$. Це виконується лише при $x=1$. Відповідь: \textbf{А}.}\fi
\par\penalty-20
\end{samepage}

\begin{samepage}
% НМТ 2026, сесія 15.06, завдання №7
\zadtask{Розв'яжіть нерівність $5^{x-1} \leqslant \dfrac{1}{25}$. \nmtyear{2026}}
\answerTable{$(-\infty; -1]$}{$(-\infty; -4]$}{$(-\infty; -3]$}{$[1; +\infty)$}{$[3; +\infty)$}
% Відповідь: А
\par\penalty-20
\end{samepage}

\begin{samepage}
% НМТ 2026, сесія 16.06, завдання №7
\zadtask{Розв'яжіть нерівність $8^{2x+9} > 8^x$. \nmtyear{2026}}
\answerTable{$(-\infty; -3)$}{$(-3; +\infty)$}{$(-9; +\infty)$}{$(9; +\infty)$}{$(-\infty; -9)$}
% Відповідь: Г
\par\penalty-20
\end{samepage}

\clearpage
\sectionTitle{Відповіді}

\noindent{\small\color{gray!80!black}Відповіді до завдань цього збірника. Для завдань НМТ-2026 узято офіційні ключі, для 2023--2025 --- обчислено.}\par\vspace{0.3cm}

\begin{multicols}{5}\noindent
\mbox{1~--- Б}\par
\mbox{2~--- В}\par
\mbox{3~--- А}\par
\mbox{4~--- Д}\par
\mbox{5~--- В}\par
\mbox{6~--- А}\par
\mbox{7~--- Д}\par
\mbox{8~--- В}\par
\mbox{9~--- А}\par
\mbox{10~--- Д}\par
\mbox{11~--- Б}\par
\mbox{12~--- Б}\par
\mbox{13~--- Б}\par
\mbox{14~--- Б}\par
\mbox{15~--- А}\par
\mbox{16~--- Г}\par
\mbox{17~--- Д}\par
\mbox{18~--- Г}\par
\mbox{19~--- В}\par
\mbox{20~--- А}\par
\mbox{21~--- В}\par
\mbox{22~--- В}\par
\mbox{23~--- Б}\par
\mbox{24~--- Г}\par
\mbox{25~--- Г}\par
\mbox{26~--- Б}\par
\mbox{27~--- Г}\par
\mbox{28~--- Д}\par
\mbox{29~--- Г}\par
\mbox{30~--- Д}\par
\mbox{31~--- Б}\par
\mbox{32~--- Г}\par
\mbox{33~--- Б}\par
\mbox{34~--- В}\par
\mbox{35~--- Д}\par
\mbox{36~--- Б}\par
\mbox{37~--- Б}\par
\mbox{38~--- Д}\par
\mbox{39~--- В}\par
\mbox{40~--- А}\par
\mbox{41~--- Д}\par
\mbox{42~--- Г}\par
\mbox{43~--- А}\par
\mbox{44~--- Б}\par
\mbox{45~--- Б}\par
\mbox{46~--- Б}\par
\mbox{47~--- Б}\par
\mbox{48~--- А}\par
\mbox{49~--- Г}\par
\mbox{50~--- А}\par
\mbox{51~--- А}\par
\mbox{52~--- В}\par
\mbox{53~--- В}\par
\mbox{54~--- Д}\par
\mbox{55~--- Б}\par
\mbox{56~--- В}\par
\mbox{57~--- А}\par
\mbox{58~--- Б}\par
\mbox{59~--- Б}\par
\mbox{60~--- Г}\par
\mbox{61~--- Д}\par
\mbox{62~--- Г}\par
\mbox{63~--- А}\par
\mbox{64~--- Г}\par
\mbox{65~--- Б}\par
\mbox{66~--- А}\par
\mbox{67~--- Г}\par
\mbox{68~--- Г}\par
\mbox{69~--- Д}\par
\mbox{70~--- Д}\par
\mbox{71~--- А}\par
\mbox{72~--- Д}\par
\mbox{73~--- А}\par
\mbox{74~--- А}\par
\mbox{75~--- В}
\end{multicols}

\end{document}
